\chapter{Discusión y conclusiones}
\label{cap:discusion}
\estadoBorrador

La discusión final permanece abierta; esta etapa identifica preguntas
sustentadas por resultados exploratorios. La propuesta predictiva anterior
se conserva únicamente como antecedente posible. Una eventual discusión de
esa línea se redactaría a partir de resultados evaluados y aprobados.
En esa posible línea futura se contrastarían las hipótesis y se identificarían
limitaciones de selección y
medición, delimitar la aplicabilidad de las alertas y separar asociación
predictiva de inferencia causal. La auditoría actual sustenta la necesidad
de distinguir modalidades, fechas y poblaciones; no demuestra todavía una
ganancia predictiva ni una anticipación útil para intervenir.

% BEGIN CONSOLIDATED 05_preguntas_exploratorias
\section{Posibles líneas de investigación derivadas de los hallazgos}

\label{sec:exploracion-preguntas}

Las siguientes líneas se derivan de los agregados observados y permanecen abiertas. Su prioridad depende de interés sustantivo, cobertura y validación, no de seleccionar un outcome predeterminado. No se estima un modelo ni se aprueba una pregunta definitiva.

\subsection{¿Cómo difiere la recepción propia y del cónyuge por sexo y escolaridad inicial?}

Resultado que motiva: T1\_2021: brecha propia de 18,59 pp en 6 años y 0,55 pp en $\geq$10 años.

Variables y rondas: K propia/cónyuge, sexo, escolaridad inicial, edad, estado conyugal. Perfiles transversales seis rondas; trayectorias condicionadas K2001--2003/2012--2015/2018--2021.

Viabilidad: Validar J1/2001, filtros de pareja, selección y antecedente laboral. Descriptiva estratificada; econométrica asociacional tras fijar estimando; predicción solo futura y opcional.

\subsection{¿En qué perfiles coexisten pensiones K y apoyos, y cómo varía la coexistencia con la edad?}

Resultado que motiva: 2021: 2\,019 K+A; 3\,558 A con negativa K; V=0,03756 y tasas de A muy diferentes por edad quinquenal.

Variables y rondas: K/A propios, edad, sexo, escolaridad, localidad; programa explícito cuando existe. Cortes seis rondas con contratos separados; recepción conjunta 2012--2015 condicionada; A2018--2021 no homologado.

Viabilidad: No llamar sin pensión a negativa K ni todo A no contributivo; separar dinero/especie requeriría fuente adicional. Descriptiva/econométrica sin causalidad por defecto.

\subsection{¿Qué continuidad e inicio de apoyo se observa entre personas con negativa K inicial?}

Resultado que motiva: Cohorte con negativa K en 2015, 60+ viva: 5\,503 personas. Inicio A observado 653/2\,740 con negativa A inicial y respuesta final conocida (23,83\%). Mismo donativo adicional, dinero/especie; 2018 desglosa hasta dos instituciones. Solo recepción agregada condicionada.. No equivale a personas sin toda pensión, ni implica que el apoyo haya causado ampliación de cobertura.

Variables y rondas: K en origen, A en ambos momentos, perfiles iniciales, seguimiento y modalidades. A2003--2012/2012--2015/2015--2018 condicionados; intervalos nominales desiguales.

Viabilidad: Necesita examinar selección de pares conocidos, retornos, muerte y sensibilidad. Descriptiva panel; modelo de transición asociacional si se define proceso de observación y estimando.

\subsection{¿Cómo se combinan recepción pensionaria y trabajo durante el envejecimiento?}

Resultado que motiva: 2021: 559/2\,760 (20,25\%) K propia entre quienes trabajan; T7 presenta heterogeneidad por edad.

Variables y rondas: Ocupación construida, edad, K propia/cónyuge, salario principal en elegibles, instituciones. Actividad en seis ondas; recepción K longitudinal solo en periodos compatibles.

Viabilidad: Ocupación no profesión; retiro/incapacidad/no trabajo no separables. Validar saltos y periodos salarial/pensionario. Descripción de perfiles y transiciones; econometría asociacional, no efecto retiro.

\subsection{¿Qué diferencias institucionales aparecen entre receptoras según escolaridad y localidad?}

Resultado que motiva: 2021: IMSS 2\,540/3\,592; instituciones múltiples no aditivas. C1/C2 usan perfiles simultáneos.

Variables y rondas: Instituciones fuente, contingencia, sexo, edad, educación, localidad, titularidad. Seis cortes por instrumento; cambios institucionales individuales requieren homologación adicional, especialmente2012/2015.

Viabilidad: Fuente no régimen y código ambiguo permanece desconocido. No asignar monto por institución cuando selección por contingencia. Descriptiva; transición institucional pendiente, no entrenar un modelo.

\subsection{¿Cómo varía la asociación de monto pensionario propio y patrimonio compartido por escolaridad?}

Resultado que motiva: 2021: Spearman 0,2643 en 3\,347 pares; T8 distribuye patrimonio simultáneamente por ramas y educación.

Variables y rondas: Patrimonio, monto K parcial, escolaridad, tipo/slot, titular y pareja. Distribuciones internas a seis rondas; cambio real no disponible hasta validar moneda/precios y unidad compartida.

Viabilidad: No vulnerabilidad por cuartil, no ingreso total y no acumulación causal. Descriptiva actual; longitudinal monetaria pendiente de deflactor y consistencia de unidad.
% END CONSOLIDATED 05_preguntas_exploratorias
